% !LW recipe = latexmk (xelatex)
\documentclass[UTF8,a4paper,10pt]{ctexart}
\usepackage[margin=14mm,top=13mm,bottom=15mm]{geometry}
\setlength{\columnsep}{5mm}
\setlength{\parindent}{1em}
\setlength{\parskip}{0pt}
\linespread{0.96}
\ctexset{
  section={format=\raggedright\bfseries\large,beforeskip=7pt,afterskip=3pt},
  subsection={format=\raggedright\bfseries\normalsize,beforeskip=5pt,afterskip=2pt}
}
\usepackage{microtype}
\usepackage{enumitem}
\setlist[itemize]{nosep}
\usepackage{xcolor}
\definecolor{attackRed}{HTML}{C03736}
\definecolor{threatYellow}{HTML}{E5B536}
\definecolor{moveBlue}{HTML}{245E91}
\definecolor{movePurple}{HTML}{7A4397}
\usepackage{chessboard}
\usepackage{twemojis}
\usepackage{newunicodechar}
\setchessboard{boardfontsize=17pt,showmover=true,moverstyle=square}

\newfontfamily\arrowlabelfont{Latin Modern Sans}
\newcommand{\ArrowLabel}[3]{%
  \begingroup\setlength{\fboxsep}{0.5pt}%
  \colorbox{#1}{\textcolor{#2}{{\arrowlabelfont\fontsize{6}{6}\selectfont #3}}}%
  \endgroup}
% Black moves: yellow arrow and label; White moves: blue arrow and label.
\newcommand{\BoardMoveArrow}[5]{%
  \pgfsetstrokecolor{#2}%
  \pgfsetlinewidth{1.1pt}%
  \pgfsetarrowsend{stealth}%
  \pgfpathmoveto{\pgfpointorigin}%
  \pgfpathlineto{\pgfpointanchor{#1}{center}}%
  \pgfusepath{stroke}%
  \pgftext[at=\pgfpointscale{0.5}{\pgfpointanchor{#1}{center}}]{%
    \ArrowLabel{#3}{#4}{#5}}}
\cbDefinePgfMoveStyle{plyone}{\BoardMoveArrow{#1}{threatYellow}{threatYellow}{black}{1}}
\cbDefinePgfMoveStyle{plytwo}{\BoardMoveArrow{#1}{moveBlue}{moveBlue}{white}{2}}
\cbDefinePgfMoveStyle{plythree}{\BoardMoveArrow{#1}{threatYellow}{threatYellow}{black}{2}}
\cbDefinePgfMoveStyle{plyfour}{\BoardMoveArrow{#1}{moveBlue}{moveBlue}{white}{3}}
\cbDefinePgfMoveStyle{whiteplyone}{\BoardMoveArrow{#1}{moveBlue}{moveBlue}{white}{1}}
\newcommand{\BoardInverseMoveArrow}[7]{%
  \pgfsetstrokecolor{#2}%
  \pgfsetlinewidth{1.1pt}%
  \pgfsetarrowsend{stealth}%
  \pgfpathmoveto{\pgfpointorigin}%
  \pgfpathlineto{\pgfpointanchor{#1}{center}}%
  \pgfusepath{stroke}%
  \pgftext[rotate=180,at=\pgfpointadd{\pgfpointscale{0.5}{\pgfpointanchor{#1}{center}}}{\pgfpoint{#6}{#7}}]{%
    \ArrowLabel{#3}{#4}{#5}}}
\cbDefinePgfMoveStyle{backrankqueen}{\BoardInverseMoveArrow{#1}{threatYellow}{threatYellow}{black}{1}{10pt}{0pt}}
\cbDefinePgfMoveStyle{backrankwhiterook}{\BoardInverseMoveArrow{#1}{moveBlue}{moveBlue}{white}{2}{0pt}{-4pt}}
\cbDefinePgfMoveStyle{backrankblackrook}{\BoardInverseMoveArrow{#1}{threatYellow}{threatYellow}{black}{2}{-8pt}{0pt}}
% Shift the two rook rays to opposite sides of the f-file so both remain visible.
\newcommand{\StackedRookArrow}[2]{%
  \pgfsetstrokecolor{attackRed}%
  \pgfsetlinewidth{1.2pt}%
  \pgfsetarrowsend{stealth}%
  \pgfpathmoveto{\pgfpointadd{\pgfpointorigin}{\pgfpoint{#2}{0pt}}}%
  \pgfpathlineto{\pgfpointadd{\pgfpointanchor{#1}{center}}{\pgfpoint{#2}{0pt}}}%
  \pgfusepath{stroke}}
\cbDefinePgfMoveStyle{rookrayleft}{\StackedRookArrow{#1}{-2pt}}
\cbDefinePgfMoveStyle{rookrayright}{\StackedRookArrow{#1}{2pt}}
\cbDefinePgfMoveStyle{queenroute}{%
  \pgfsetstrokecolor{moveBlue}%
  \pgfsetlinewidth{1.1pt}%
  \pgfsetarrowsend{stealth}%
  \pgfpathmoveto{\pgfpointorigin}%
  \pgfpathlineto{\pgfpointanchor{#1}{center}}%
  \pgfusepath{stroke}}
\cbDefinePgfRegionStyle{dashedredbox}{%
  \pgfsetstrokecolor{attackRed}%
  \pgfsetlinewidth{1pt}%
  \pgfsetdash{{2pt}{1.5pt}}{0pt}%
  \pgfpathrectanglecorners{\pgfpointxy{-0.44}{0.44}}%
    {\pgfpointadd{\pgfpointanchor{#1}{center}}{\pgfpointxy{0.44}{-0.44}}}%
  \pgfusepath{stroke}}
% An unnumbered dashed arrow shows an attack opened by a move.
\cbDefinePgfMoveStyle{attackray}{%
  \pgfsetstrokecolor{attackRed}%
  \pgfsetlinewidth{1pt}%
  \pgfsetdash{{2pt}{1.5pt}}{0pt}%
  \pgfsetarrowsend{stealth}%
  \pgfpathmoveto{\pgfpointorigin}%
  \pgfpathlineto{\pgfpointanchor{#1}{center}}%
  \pgfusepath{stroke}}
\newcommand{\Position}[3][]{%
  \par\noindent\makebox[\columnwidth][c]{\chessboard[setfen={#2},#1]}\par\nobreak
  \noindent\makebox[\columnwidth][c]{{\footnotesize\bfseries Fig.~\thesubsection\quad #3}}\par%
}

\begin{document}
\pagestyle{plain}
\twocolumn[{%
  \begin{center}
    {\Large\bfseries Checkmate Patterns\par}
    \vspace{2pt}
    {\small Tactical Motifs \quad\textbar\quad \the\year/\the\month/\the\day\par}
  \end{center}
  \vspace{-7pt}
}]

%%%%%%%%%%%%%%%%%%%%%%%%%%%%%%%%%%%%%%%%%%%%%%%%%%%%%%%%%%%
% Fool's Mate
%%%%%%%%%%%%%%%%%%%%%%%%%%%%%%%%%%%%%%%%%%%%%%%%%%%%%%%%%%%
% 愚人杀法的核心：f 兵缺失，后和象借王翼斜线攻击；弃后引开并消除 h 兵。
\section[Fool's Mate]{\texttwemoji{clown face} Fool's Mate}
\subsection{Queen Sacrifice on g6}

With the f7-pawn gone, the h5--e8 diagonal reaches the black king.

\Position[
  backstyle=color,color=attackRed!42,backfields={g6,f7,e8},
  markstyle=circle,color=movePurple,markfields={h7},
  markstyle=attackray,markmove=h5-e8,
  markstyle=plyone,markmove=g7-g6,
  markstyle=plytwo,markmove=h5-g6,
  markstyle=plythree,markmove=h7-g6,
  markstyle=plyfour,markmove=d3-g6
]{rnbqkb1r/ppppp1pp/8/3n3Q/8/3B4/PPP2PPP/R1B1K1NR b - - 0 1}{1...g6 2.Qxg6+ hxg6 3.Bxg6\#}

%如果出了后兵却仍然挡着后（罕见情况），也可以逼出f兵。
\subsection{Knight Sacrifice on e6}

If the queen's pawn has advanced but the queen is still blocked (a rare case), the f-pawn can also be drawn out.

Earlier, the knight was on g5. Against:
\begin{itemize}
  \item ...f6: Ne6 can get the trapped queen!
  \item ...h6: on above case.
\end{itemize}

\Position[
  backstyle=color,color=attackRed!42,backfields={d8,f8},
  markstyle=circle,color=movePurple,markfields={h6},
  markstyle=attackray,markmove=h5-e8,
  markstyle=plyone,markmove=f7-e6,
  markstyle=plytwo,markmove=d1-h5,
  markstyle=plythree,markmove=g7-g6,
  markstyle=plyfour,markmove=h5-g6
]{r1bqkbnr/pppnppp1/4N2p/8/4p3/2N5/PPPP1PPP/R1BQKB1R b - - 0 1}{1...fxe6 2.Qh5+ g6 3.Qxg6\#}

%%%%%%%%%%%%%%%%%%%%%%%%%%%%%%%%%%%%%%%%%%%%%%%%%%%%%%%%%%%
% Back Rank Mate
%%%%%%%%%%%%%%%%%%%%%%%%%%%%%%%%%%%%%%%%%%%%%%%%%%%%%%%%%%%

% The key is to seize the back rank. The amount of force depends on how many pieces are stacked and connected; however many there are, they are useless if scattered.
% Do not let the opponent's rook and queen stack up; if they do, the back rank cannot be defended.
% The stacking here is unusual: starting from the attack point on the back rank, stack the rook and queen along a file, and the queen along a diagonal. So it is not really "stacking"; the core is concentrating firepower.
  % This algorithm is complicated; the opponent's force must also be calculated for each square.
  % Calculate it this way: for each back-rank square, add +1/-1 for occupation and attack. Finding or creating a square with a score greater than 0 means mate is possible.
% Do not let the opponent's rook and queen reach the king's side; the king can block them.
\section{Back-Rank Mate}


% 关键点是集中火力到格子，以及引离保护底线的车后。
The key is to \textcolor{red}{concentrate firepower} and \textcolor{red}{draw away the rook and queen protecting the back rank}.

\subsection{Bobby Fisher's Game}

Qf2 not only has a chance to draw away the f1 rook, but also blocks the opponent's attack on f7.

\begin{figure}[t]
\Position[
  inverse=true,
  backstyle=color,color=attackRed!42,backfields={f7,e1},
  markstyle=attackray,markmove=f2-e1,
  markstyle=backrankqueen,markmove=d4-f2,
  markstyle=backrankwhiterook,markmove=f1-g1,
  markstyle=backrankblackrook,markmove=e8-e1
]{4rbk1/3Q1ppp/3p4/3P4/3q4/B7/P5PP/5R1K b - - 0 1}{1...Qf2 2.Rg1 Re1}
\end{figure}

\subsection{Overload Defender Rook}

f8 rook would suffer from tripple pressure now.

\begin{figure}[t]
\Position[
  markstyle=rookrayleft,markmove=f7-f8,
  markstyle=rookrayright,markmove=f1-f8,
  markstyle=attackray,markmove=e7-e8,
  markstyle=attackray,markmove=e7-f8,
  markstyle=whiteplyone,markmove=c7-e7
]{4rrk1/2Q2Rpp/2p5/p7/3q4/8/2P3PP/5R1K w - - 0 1}{1. Qe7}
\end{figure}

%%%%%%%%%%%%%%%%%%%%%%%%%%%%%%%%%%%%%%%%%%%%%%%%%%%%%%%%%%%
% Kiss of Death TODO
%%%%%%%%%%%%%%%%%%%%%%%%%%%%%%%%%%%%%%%%%%%%%%%%%%%%%%%%%%%

%%%%%%%%%%%%%%%%%%%%%%%%%%%%%%%%%%%%%%%%%%%%%%%%%%%%%%%%%%%
% Damiano's Mate
%%%%%%%%%%%%%%%%%%%%%%%%%%%%%%%%%%%%%%%%%%%%%%%%%%%%%%%%%%%

\section{\texttwemoji{alembic} Damiano's Mate}

% 这个模式很简单，关键在于g6兵对h7后的支持，以及f8挡住王的去路（一般是易位后的车）。
% 为了达成这个局面，我们往往需要创造h开放线，并常伴随着h车弃子来给后开路。
This pattern is quite straightforward; the key lies in the support provided by the g6 pawn to the h7 square, combined with the f8 square being blocked (typically by a rook from castling).

To reach this position, one usually needs to open h-file, which frequently involves sacrificing h-file rook, to clear a path for the queen.

\begin{figure}[t]
\noindent\makebox[\columnwidth][c]{\chessboard[
  setfen={3r1rk1/1p3pp1/1qp3P1/2nnp3/p7/3B1P2/PPP3P1/1K1QR2R w - - 0 1},
  showmover=false,
  markstyle=queenroute,markmove=d1-h1,
  markstyle=queenroute,markmove=h1-h7,
  markstyle=attackray,markmove=h7-g8,
  markstyle=attackray,markmove=h7-h8,
]}\par
\noindent\makebox[\columnwidth][c]{{\footnotesize\bfseries Fig.~\thesection}}\par
\end{figure}

%%%%%%%%%%%%%%%%%%%%%%%%%%%%%%%%%%%%%%%%%%%%%%%%%%%%%%%%%%%
% Opera Mate
%%%%%%%%%%%%%%%%%%%%%%%%%%%%%%%%%%%%%%%%%%%%%%%%%%%%%%%%%%%


\end{document}
